\documentclass[11pt]{article}
\usepackage[margin=1in]{geometry}
\usepackage{amsmath, amssymb}
\usepackage{booktabs}
\usepackage{multirow}
\usepackage{parskip}
\usepackage{titlesec}
\titlespacing*{\section}{0pt}{6pt}{3pt}
\titlespacing*{\subsection}{0pt}{4pt}{2pt}

\title{\textbf{Predict-then-Optimize for Stochastic DC-OPF}\\[4pt]
\large A Two-Stage Stochastic Programming Framework}
\author{}
\date{}

\begin{document}
\maketitle
\vspace{-1.5em}

\section{Overview}

We consider the problem of committing generator dispatch for the IEEE 118-bus system under uncertain nodal demand. Let $N$, $G$, $E$, $S$ denote the sets of buses, generators, branches, and scenarios respectively. The demand at each bus $n$ is a random variable $P^d_{n}$ whose distribution we approximate via a finite scenario set $\{P^d_{n,s}\}_{s \in S}$. We compare three methods for constructing this scenario set.

\section{Three-Phase Methodology}

\textbf{Phase 1 — Na\"ive Lognormal (Baseline).} System-level demand $P^d$ is assumed to follow a lognormal distribution fitted to historical CAISO data. Node-level scenarios are obtained by scaling a fixed nodal allocation vector: $P^d_{n,s} = w_n \cdot P^d_s$, where $w_n$ is the historical average share of bus $n$.

\textbf{Phase 2 — Predict-then-Optimize (PGBM).} A 17-quantile gradient boosting model (PGBM) is trained on the stationary residual of CAISO load after STL decomposition at periods of 96 and 672 intervals. At each evaluation timestep $t$, the model produces a predictive distribution over load residuals; scenarios are reconstructed in MW space by adding back the trend and seasonal components. The resulting $|S| = 100$ scenarios $\{P^d_{n,s}\}$ replace the lognormal draws.

\textbf{Phase 3 — Smart PtO (forthcoming).} End-to-end differentiable optimization that back-propagates the operational cost through the LP solver into the forecasting model's training loss, directly aligning forecast quality with decision quality.

\section{Two-Stage Stochastic DC-OPF}

Given scenarios $\{P^d_{n,s}\}_{s\in S}$, we solve:

\begin{equation*}
\min_{\tilde{P},\,\theta_1,\,P_1,\,\theta_2,\,P_2,\,G^+,\,G^-}
\;\sum_{g\in G}\!\bigl(C_{g,1}\tilde{P}_g + C_{g,0}\bigr)
+ \frac{1}{|S|}\sum_{s\in S}\!\left[
  \sum_{g\in G} C_{g,1}\bar{\alpha}_g\delta_s
  + \sum_{n\in N}\!\bigl(C^{\text{voll}}G^+_{n,s} + C^{\text{curt}}G^-_{n,s}\bigr)
\right]
\end{equation*}

\textbf{Stage 1 — Commitment} (decisions taken before demand is observed):
\begin{align*}
&\theta_{\text{ref},1} = 0, \qquad
 P_{e,1} = B_e\!\left(\theta_{o(e),1} - \theta_{d(e),1}\right)\;\forall e,\\
&\sum_{g\in G_n}\!\tilde{P}_g - \hat{P}^d_n
 = \!\!\sum_{e\in E^{\text{from}}_n}\!\!P_{e,1} - \!\!\sum_{e\in E^{\text{to}}_n}\!\!P_{e,1}\;\forall n,\qquad
 P^{\min}_g \le \tilde{P}_g \le P^{\max}_g\;\forall g,\qquad
 |P_{e,1}| \le R_e\;\forall e.
\end{align*}

Here $\hat{P}^d_n = \mathbb{E}_s[P^d_{n,s}]$ is the point forecast (lognormal mean or PGBM median, depending on phase).

\textbf{Stage 2 — Recourse} (after scenario $s$ is revealed, generation adjusts via affine policy $\tilde{P}_g + \bar{\alpha}_g\delta_s$):
\begin{align*}
&\delta_s = \textstyle\sum_n P^d_{n,s} - \sum_n \hat{P}^d_n,\qquad
 \bar{\alpha}_g = P^{\max}_g \big/ \textstyle\sum_{g'}P^{\max}_{g'}\quad\text{(fixed participation factors)},\\
&\theta_{\text{ref},s,2} = 0,\qquad
 P_{e,s,2} = B_e\!\left(\theta_{o(e),s,2} - \theta_{d(e),s,2}\right)\;\forall e,s,\\
&P^{\min}_g \le \tilde{P}_g + \bar{\alpha}_g\delta_s \le P^{\max}_g\;\forall g,s,\\
&\sum_{g\in G_n}\!\!\bigl(\tilde{P}_g+\bar{\alpha}_g\delta_s\bigr) - \bigl(P^d_{n,s} - G^+_{n,s}\bigr) - G^-_{n,s}
 = \!\!\sum_{e\in E^{\text{from}}_n}\!\!P_{e,s,2} - \!\!\sum_{e\in E^{\text{to}}_n}\!\!P_{e,s,2}\;\forall n,s,\\
&|P_{e,s,2}| \le R_e\;\forall e,s,\qquad
 0 \le G^+_{n,s} \le P^d_{n,s}\;\forall n,s,\qquad
 G^-_{n,s}\ge 0\;\forall n,s.
\end{align*}

Slack variables $G^+_{n,s}$ (load shed) and $G^-_{n,s}$ (curtailment) are penalised at $C^{\text{voll}}=\$10{,}000$/MWh and $C^{\text{curt}}=\$1{,}000$/MWh respectively. The LP is solved with HiGHS via \texttt{scipy.optimize.linprog}, with $|N|=118$, $|G|=54$, $|E|=186$, $|S|=100$, yielding ${\approx}54{,}000$ decision variables.

\section{Preliminary Results (Phase 2, $T=152$ of 200 timesteps)}

In Phase 2, the two-stage stochastic DC-OPF is solved identically to Phase 1; the only difference is that the scenario set $\{P^d_{n,s}\}$ is generated by PGBM rather than the na\"{i}ve lognormal. PGBM is trained independently to minimize quantile (pinball) loss, so the forecaster and optimizer remain decoupled. Phase 3 will close this gap by training the forecaster end-to-end with the optimizer.

\begin{center}
\begin{tabular}{llr}
\toprule
Category & Metric & Value \\
\midrule
\multirow{4}{*}{Forecast quality}
 & MAE & 227.8 MW \\
 & Bias & $-37.0$ MW \\
 & Avg.\ scenario std.\ $\hat{\sigma}_s$ & 351.7 MW \\
 & Timesteps solved / attempted & 152 / 200 \\
\midrule
\multirow{3}{*}{Operational cost}
 & $\mathbb{E}[\text{SP expected cost}]$ & \$109,836/hr \\
 & $\mathbb{E}[\text{realized cost}]$ & \$101,664/hr \\
 & Calibration gap (realized $-$ expected) & $-$\$8,171/hr \\
\midrule
\multirow{4}{*}{Realized performance}
 & Avg.\ committed generation $\sum_g \tilde{P}_g$ & 3,816.6 MW \\
 & Avg.\ actual demand $\sum_n P^d_n$ & 3,853.6 MW \\
 & Avg.\ load shed $G^+$ & 0.00 MW \\
 & Avg.\ line-limit violation & 0.28 MW \\
\bottomrule
\end{tabular}
\end{center}

The negative calibration gap indicates that the SP slightly over-provisions generation: committed dispatch exceeds realized demand on average by $\sim$37 MW, leaving unused headroom. The 48 failed solves arise at high-uncertainty timesteps where the scenario spread $\delta_s^{\max} - \delta_s^{\min}$ exceeds individual generator capacity ranges, creating an infeasible affine recourse problem; these are handled by falling back to the original $[P^{\min}_g, P^{\max}_g]$ bounds. Phase 3 is expected to reduce both the calibration gap and infeasibility rate by aligning the forecast distribution with operational requirements.

\end{document}
